\mychapter{0}{Résumé}

\begin{flushleft}

Le présent travail a été réalisé dans le cadre du projet de fin d’études pour l’obtention du diplôme d’ingénieur en informatique. Ce projet porte sur le développement de TSS, une plateforme intelligente dédiée à la gestion et au suivi des services techniques ainsi qu’à la gestion des tickets d’incidents.  

L’objectif principal de cette application est d’optimiser le processus de traitement des demandes techniques en facilitant la communication entre les différents acteurs, notamment les coordinateurs, les techniciens et le client. La plateforme permet la création, le suivi et la gestion des tickets, tout en assurant une meilleure organisation des interventions et une amélioration de la qualité des services.  

Le système intègre plusieurs fonctionnalités importantes, telles que la gestion des utilisateurs selon leurs rôles, la gestion des tickets et des équipements, le suivi des interventions techniques, la génération de statistiques et de tableaux de bord, ainsi qu’un système de notifications en temps réel. Une interface web moderne et responsive a également été développée afin d’offrir une expérience utilisateur fluide et intuitive.  

Mots clés : plateforme web, gestion des tickets, maintenance, suivi des interventions, notifications, tableau de bord, application intelligente.\\[1pt]

\medskip
\medskip
\medskip

\Huge{Abstract}

\medskip
\medskip
\medskip

This work was carried out as part of the final-year project for obtaining the Computer Engineering degree. The project focuses on the development of TSS, an intelligent platform dedicated to technical service management and incident ticket handling.  

The main objective of this application is to optimize the processing of technical requests by improving communication between different actors, including coordinators, technicians, and helpdesk agents. The platform allows the creation, tracking, and management of tickets while ensuring better organization of technical interventions and improving service quality.  

The system integrates several important features such as user management based on roles, ticket and equipment management, monitoring of technical interventions, generation of statistics and dashboards, and a real-time notification system. A modern and responsive web interface was also developed to provide a smooth and intuitive user experience.  

Keywords: web platform, ticket management, maintenance, intervention tracking, notifications, dashboard, intelligent application.

\end{flushleft}



